% Source-aligned English for printed p.205.
\clearpage
\SourcePage{205}

\LayerTarget{tohoku:proof:corollary-5.3.1:2:p1}
\noindent We have $I_2^{p,q}=0\ \mathrm{mod}\ {\mathbf C}(P)$ for $q>0$, since
by Theorem 5.3.1 $\mathbf H^q(G,A)$ is a sheaf annihilated by $n$; hence
$H^p(Y,\mathbf H^q(G,A))$ is annihilated by $n$ for every $p$, and is
therefore zero modulo ${\mathbf C}(P)$. Corollary 2 follows at once.

\LayerTarget{tohoku:application:fixed-point-free-case-5.3}
If all the $n_x$ are equal to 1, that is, if every $g\in G$, $g\ne e$, has no
fixed point (we then simply say that $G$ \textit{acts without fixed points}),
we obtain $\mathbf H^n(G,A)=0$ for $n>0$, and hence
$H^*(X,G,A)=H^*(Y,A^G)$. This is also evident a priori from the isomorphism,
noted at the end of 5.1, between the category of sheaves on $Y$ and the
category of $G$-sheaves on $X$ when $G$ is a discrete group of homeomorphisms
without fixed points. We therefore obtain in this case the classical
statement:

\LayerTarget{tohoku:corollary:5.3.1:3}
\textsc{Corollary 3.} \textit{If $G$ acts ``without fixed points'',
then $H^*(Y,A^G)$ is the abutment of a cohomological spectral sequence whose
initial term is $H^p(G,H^q(X,A))$.}

\LayerTarget{tohoku:remark:spectral-sequence-construction-5.3}
\textsc{Remark.} By definition, the second spectral sequence of
Theorem 5.2.1 is obtained by taking a resolution $C(A)$ of $A$ by
$\Gamma_X$-acyclic sheaves and taking the second spectral sequence of the
functor $\Gamma^G$ relative to the complex $\Gamma_X(C(A))$. In general, the
first spectral sequence of $\Gamma^G$ relative to this complex is clearly not
identified with the first spectral sequence of Theorem 5.2.1 (take an
injective resolution of $A$!). It is, however, so identified if one takes for
$C(A)$ the \textit{canonical resolution} of $A$ (see 3.3), as follows by an
argument analogous to that of Proposition 4.3.2. It is also so identified if
$X$ and $Y$ are paracompact and one takes for $C(A)$ a Cartan resolution of
$A$, i.e. $C(A)=A\otimes_ZC$, where ${\mathbf C}$ is a ``fundamental sheaf''
in the sense of Cartan--Leray, as can be shown by rather different arguments
(we shall not need this fact). Consequently spectral sequence I is precisely
the spectral sequence of [10], whose initial term had been left unspecified.

\LayerTarget{tohoku:section:5.4}
\bigskip
\textbf{5.4. Transformation of the first spectral sequence.} We
continue to suppose that condition (D) of the preceding section holds
(although this is not absolutely essential).
\LayerTarget{tohoku:setup:transformation-first-spectral-sequence}
Let $Y_0$ be a closed subset of $Y$ containing the supports of the
$\mathbf H^n(G,A)$ for $n>0$ (for example, by Theorem 5.3.1 it is enough that
$Y_0$ contain the set of $y\in Y$ such that $G_y\ne e$). Let
$X_0=f^{-1}(Y_0)$, let $V$ be the complement of $Y_0$, and let $U$ be the
complement of $X_0$. Consider the exact sequences

\LayerTarget{tohoku:sequence:5.4-support-exact-sequences}
\[
\begin{gathered}
0\to A_U\to A\to A_{X_0}\to0,\\
0\to(A^G)_V\to A^G\to
  (A^G)_{Y_0}\to0.
\end{gathered}
\]
Since manifestly $(A_U)^G=(A^G)_V$, we obtain a commutative diagram of
cohomology sequences:

\LayerTarget{tohoku:diagram:5.4.1}
\[
\tag{5.4.1}
\begin{array}{c@{\;}c@{\;}c@{\;}c@{\;}c@{\;}c@{\;}c}
\mathinner{.\,.\,.\,.}H^n(Y,(A^G)_V)&\to&H^n(Y,A^G)&\to&H^n(Y,(A^G)_{Y_0})&\to&
H^{n+1}(Y,(A^G)_V)\!\to\!\mathinner{.\,.\,.\,.}\\[-0.2em]
\downarrow\wr&&\downarrow&&\downarrow&&\downarrow\wr\\[0.45em]
\mathinner{.\,.\,.\,.}H^n(X;G,A_U)&\to&H^n(X:G,A)&\to&H^n(X;G,A_{X_0})&\to&
H^{n+1}(X;G,A_U)\!\to\!\mathinner{.\,.\,.\,.}
\end{array}
\]
\LayerTarget{tohoku:proof:transformation-first-spectral-sequence:p1}
It follows from the hypothesis on $Y_0$ that $\mathbf H^n(G,A_U)=0$ for
$n>0$, since this is clearly the case on the open set $V$, but it is also true
on $Y_0$, because by

\TohokuContinuation{206}
